\documentclass[./unravbook.tex]{subfiles}

\usepackage{parskip}
\setlength{\parskip}{10pt}
\setstretch{1.15}

\setenotez{
  list-heading = {\chapter*{#1}\markboth{#1}{#1}}
}

\begin{document}

\setcounter{chapter}{2}
\setcounter{endnote}{0}

\chapter{Criminalized} %% {{{
    %% end header ----------------------------
    

In the late summer of 2001, over a weeklong period ending a few days before the 9/11 attacks in the United States, the UN held its third \enquote{World Conference Against Racism} in Durban, South Africa, known (to distinguish it from a 2009 conference held to review its progress) as \enquote{Durban I.}\footnote{Report of the World Conference against Racism, Racial Discrimination, Xenophobia and Related Intolerance Durban, 31 August -- 8 September 2001, A/CONF.189/12 (Part III) \url{https://digitallibrary.un.org/record/451954/files/A_CONF.189_12\%28PartIII\%29-EN.pdf}.} Attendees came expecting ``a neutral platform'' in which they would hear about and discuss \enquote{sensitive issues that world leaders were grappling with} relating to ``the social construction of race and citizenship; the socio-economic and political forces that drive racism and inequalities; organized responses to cultural diversity; and the impact of various types of public policies on race relations.''\footnote{Yusuf Bangura and Rodolpho Stevenhagen (eds.), \textit{Racism and Public Policy (2005)} Palgrave Macmillan, p. ix.} Meetings and presentations on the conference agenda did not include discussions relating specifically to the ongoing challenges between Israel and Palestine --- that is, the by then exhaustively examined ``Palestine Problem''\footnote{E.g., Walid Khalidi, ``The Palestine Problem: An Overview,'' Autumn 1991, \textit{Journal of Palestine Studies}, 21(1), 5--16. \url{https://doi.org/10.2307/2537362}.} --- nor were they expected to cover what did, or did not, constitute apartheid. That concept had already been codified into international law for a third time when the UNGA voted a few years earlier to ratify the Rome Statute --- which established it as a crime against humanity that could be prosecuted by the then newly established International Criminal Court --- in 1998.\footnote{Durban I took place between the time that Rome Statute had been ratified in 1998 and was ``entered into force'' in 2002. Rome Statute  - Core Legal Text. International Criminal Court. \url{https://www.icc-cpi.int/publications/core-legal-texts/rome-statute-international-criminal-court}.} However, the third WCAR is now often remembered today precisely because of its effect on the \enquote{Palestine problem}, due to what occurred at a large, boisterous NGO forum held during the week preceding it. 

Conducted in a nearby facility and ``encouraged and facilitated by the UN,''\footnote{McDougall 2002.} the NGO forum published its own concluding \enquote{Declaration} --- in a layout similar to that of an official UNGA resolution --- in which Israel's policies vis-à-vis Palestinians were featured in a separate section (\P\P 160--165), although also covered elsewhere.\footnote{WCAR NGO FORUM DECLARATION, 03 September 2001. \url{https://www.adalah.org/uploads/oldfiles/eng/intladvocacy/ngoforumdecl.htm}. There does not appear to be an authoritative version of this declaration, but the version cited here --- from the website of Adalah, the Legal Center for Arab Minority Rights --- is identical to the only other version found online, published by the Asia-Pacific NGO Movement for WCAR. \url{https://www.hurights.or.jp/wcar/E/ngofinaldc.htm}.} Declaring that \enquote{this alien domination and subjugation with the denial of territorial integrity amounts to colonialism, which denies the fundamental rights of self-determination, independence and freedom of Palestinians} (\P 161) and pronouncing a need to ``end the Israeli racist system including its own brand of apartheid'' (\P 98), the statement also recognized that the \enquote{targeted victims of Israel's brand of apartheid and ethnic cleansing methods have been in particular children, women and refugees and [it must] condemn the disproportionate numbers of children and women killed and injured in military shooting and bombing attacks} (\P 161). As one participant noted, such pronouncements made the forum a success because they deepened Israel's diplomatic isolation even while the WCAR itself failed to translate that isolation into official action:\footnote{Hassan Jabareen, ``On the Problems of Arab Identity: The example of Durban,'' October 4, 2001. \textit{Al-Hayat Newspaper (London)}. \url{https://www.adalah.org/uploads/oldfiles/upfiles/2011/Hassan\%27s\%20Al\%20Hayat\%20Article\%20about\%20Durban\%20October\%202001\%20English.pdf}.}

\begin{adjustwidth}{.65cm}{.5cm}
   \begin{singlespace}
    In Durban, we were part of the Arab group. During the first week, we worked within the forum of NGOs and civil society organizations. The governmental conference was in the second week. The Arab group managed to get everything it asked for and suggested included in both the NGO forum and the final document it issued \ldots. The wording and articles \ldots went so far as to describe Israel and the Palestinian case as a case of `genocide' and `ethnic cleansing,' Zionism as a racist movement, Israel as the last apartheid regime in the world, etc. Thus the media was drawn to the Palestinian issue to the point that the Palestinian cause became the main issue at Durban. In the second week, in the governmental document, there was serious back-tracking from the first document to the extent that there was no condemnation of Israel's policy, even on the issue of the Occupation, or any demand for the resolutions of the UN itself concerning the Occupation to be implemented.
  \end{singlespace}
\end{adjustwidth}

Israel, which knew what to expect based on the agenda of an earlier regional ``preparatory committee'' meeting to which it had not been invited, attended Durban I but like the United States, it withdrew once it saw the NGO forum's developments unfold without apparent concern by other member-state delegations. Other attendees were also generally aware that the NGO forum was being used ``to demonize Israel \ldots casting doubt on whether there was a place for it in the community of nations,''\footnote{Schoenberg H. O. (2002). Demonization in Durban: The World Conference Against Racism \textit{American Jewish Yearbook 2002}, p. 85. \url{https://www.staff.city.ac.uk/p.willetts/NGOS/WCAR/SCHOENBG.PDF}.} but this was dismissed afterward by a former UN official because ``only a minority of groups was guilty of what should be considered hate speech.'' As for the American withdrawal, while unfortunate, this was attributed to ``an aggressive lobby from sectors within the Jewish-American community.''\footnote{McDougall 2002, pp. 145-146} Numerous NGOs from Central and Eastern Europe protested  the forum's statement, with one noting that it contained ``deliberate distortions''  and was ``extremely intolerant, disrespectful and contrary to the very spirit'' of the WCAR.\footnote{European Roma Rights Centre, \enquote{ERRC Dissociates Itself from WCAR NGO Forum}, September 5, 2001. \url{https://www.errc.org/press-releases/errc-dissociates-itself-from-wcar-ngo-forum}.} 

For them and many others, watching a gathering of member-state delegations from around the world be overshadowed by NGOs \enquote{on the sidelines} who had exhorted them to address ``Israeli racism''\footnote{E.g., Hanan Ashrawi To World Conference Against Racism (31 Aug. 2001). \textit{Scoop World Independent News}. \url{https://www.scoop.co.nz/stories/WO0108/S00115/hanan-ashrawi-to-world-conference-against-racism.htm}. In Ms. Ashawi's words, ``As the world watches, Israel has succeeded in evolving and imposing another grand deception in the form of an official spin that not only dehumanizes and demonizes the Palestinians, but also as an attempt at blaming the victim and resuscitating labels that represent us as subhuman species, and genetically violent terrorists, hence undeserving of any human treatment''.} 
 was a sobering moment.\footnote{Lahav Harkov, \enquote{20 years since Durban: Most sickening display of Jew-hate since Nazis,} September 24, 2021, The Jerusalem Post. \url{https://www.jpost.com/diaspora/antisemitism/20-years-since-durban-most-sickening-display-of-jew-hate-since-nazis-680016}.} But to many UN officials, Durban I was a success and the kerfuffle over anti-Israel agitation was simply a distracting irritant to \enquote{an exhausting, exhilarating and difficult nine days and nights in Durban that ended with a breakthrough agreement.}\footnote{Durban 2001 - United Against Racism. Newsletter of the World Conference against Racism Secretariat. Issue 6, October 2021.\textit{United Nations Office of the High Commissioner for Human Rights}. \url{https://www.hurights.or.jp/wcar/E/doc/wcrnewsletter6_en.pdf}.} A darker story nonetheless emerges from what happened both before and after Durban I. In February 2001, as some regional delegates prepared for the discussions there, they considered how best to showcase the conclusions of the Report of the Asia Preparatory Committee Conference.\footnote{Steinberg, G.M. (2006). Soft Powers Play Hardball: NGOs Wage War against Israel. Israel Affairs, Vol. 12:4, pp. 748-768, p. 756.} As developed at their pre-conference meeting --- held in Tehran and without Israel being offered the opportunity to participate --- their consensus statement called for ``the cessation of all the practices of racial discrimination to which the Palestinians and the other inhabitants of the Arab territories occupied by Israel are subjected.''\footnote{Report of the Asian Preparatory Meeting (Tehran, 19-21 February 2001), (A/CONF.189/PC.2/9), para. 21. \url{https://docs.un.org/en/A/CONF.189/PC.2/9}.} Afterward, NGOs that embraced the statement and participated in the WCAR NGO forum launched what came to be known as the \enquote{Durban Strategy,}\footnote{Gerald M. Steinberg, \enquote{Soft Powers Play Hardball: NGOs Wage War Against Israel,} Israel Affairs 12(4), October 2006, 753.} geared toward accelerating global outrage against Israel through the same boycott tactics deployed in preceding decades against South Africa.\footnote{Statement: Marking 20 Years of the NGO ``Durban Strategy'' (09 Sept. 2021). NGO Monitor. \url{https://ngo-monitor.org/press-releases/marking-20-years-of-the-ngo-durban-strategy/}.} The words of Canadian legal scholar Irwin Cotler captured the earthquake but not its aftershocks:\footnote{NGO Monitor Event: 20 Years of Hijacking Human Rights - The Lasting Impact of Durban (n.d.). \textit{NGO Monitor.}
\url{https://ngo-monitor.org/live/}.}

\begin{adjustwidth}{.65cm}{.5cm}
   \begin{singlespace}
    What happened at Durban was truly Orwellian. \ldots. That which was to be the first international human rights conference of the 21st century \ldots turned into a conference which indicted Israel as the meta-human rights violator. A conference that was to commemorate the dismantling of South Africa as an apartheid state turned into a conference calling for the dismantling of Israel as an apartheid state.
    \end{singlespace}
\end{adjustwidth}

This turn of events was neither spontaneous nor left to happenstance. Rather, with the recently ratified and soon-to-be entered into force Rome Statute, Israel's enemies were now galvanized to make apartheid the iconic centerpiece of a far-reaching demonization strategy, reserved solely for the world's only Jewish-majority state and intended to associate it with the worst memories of official, state-mandated racism. This depiction was designed to not only make modern Israel synonymous with South Africa's apartheid era, but it also delivered an ipso facto guilty verdict on a fully-formed --- that is, with enforcement mechanisms, which were previously lacking --- ``crime against humanity.'' The key sponsor of this project was Iran, which had used much of the 1990s to recruit, train and deploy proxy forces against Israel as well as strengthen its military,\footnote{Bill Gertz, \enquote{Iran's Regional Powehouse,} Air Force Magazine, June 1996. \url{https://www.airandspaceforces.com/PDF/MagazineArchive/Documents/1996/June\%201996/0696iran.pdf}.} to counter what it viewed as the threat of a ``political negotiated process in the Middle East between Israel, its Arab neighbors and the Palestinians, inaugurated at the Madrid Conference in October 1991.''\footnote{Karmon, E. (8 September 1998). Iran's Policy on Terrorism in the 1990s. Internatoinal Institute for Counter-Terrorism. \url{https://ict.org.il/irans-policy-on-terrorism-in-the-1990s/}. (``Iran immediately perceived the threat to its ideological and strategic interests: the recognition of Israel's legitimacy as a state in the Middle East and the consequent strategic, political and economic advantages to the Jewish state; the consolidation of the moderate Arab regimes and at the same time the threat to the radical Islamic movements, such as the Hizballah, if Syria joined the process in earnest; the isolation of Iran on the regional level and its defeat on the ideological ground.'')}  

Having planted the seed for a resurgent front against Israel at Durban I, Iran sought to nurture its resulting fruit further at ``Durban II'' in 2009, a conference held in Geneva as a follow-up to discuss the UNGA's progress since the 2001 WCAR. With the recently formed (2006) UN Human Rights Council (HRC) now at the helm,\footnote{Welcome to the Human Rights Council, n.d., \url{https://www.ohchr.org/en/hr-bodies/hrc/about-council}.} Durban II (which actually took place inside the HRC's offices) was chaired by Libya and vice-chaired by Cuba, with Iran having a seat on the organizing committee.\footnote{Andy Levy-Ajzenkopf, ``Canada Says No to Durban II,'' January 31, 2008. The Canadian Jewish News. \url{https://thecjn.ca/news/canada-says-no-durban-ii/}.} The worm of Israel hatred was in fact now prepared to start devouring the fruit, resulting from a near decade of coordinated agitation, beginning with a \enquote{fierce outburst} by Iranian president Mahmoud Ahmadinejad at the opening plenary session.\footnote{AFP Newswire, \enquote{Western leaders recoil in horror at Ahmadinejad's \enquote{racist} speech,} April 20, 2009
https://timesofmalta.com/article/western-leaders-recoil-in-horror-at-ahmadinejads-racist-speech.253725.} This was promptly downplayed by one European NGO ``as a concession to hardliners who wanted [him] to recant his recent semi-favourable responses'' to the recent ``outreach efforts'' of the new American president, Barack Obama.\footnote{Daniel Jeffrey, ``Iran \& America: The Clenched Fist Grasping an Olive Branch,'' May 13, 2009. \textit{Royal United Services Institute}. \url{https://www.rusi.org/explore-our-research/publications/commentary/iran-america-clenched-fist-grasping-olive-branch}.} But while Ahmadinejad's speech was also dismissed by the HRC high commissioner as simply ``partisan political rhetoric,''\footnote{Navi Pillay, ``Durban Anti-racism Conference Outcome,'' Aprit 22, 2009), Office of the High Commissioner for Human Rights. \url{https://www.ohchr.org/en/opinion-editorial/2009/10/opinion-piece-durban-anti-racism-conference-outcome-navi-pillay-united}.} his words left no real reason to doubt Iran's future course of action. It was a siren call to its growing list of proxy forces in the region, and their Western supporters, to join arms in the long-term mission of bringing about Israel's eventual downfall:\footnote{
Ahmadinejad speech: full text (21 Apr. 2009) \textit{BBC News}. 
\url{http://news.bbc.co.uk/2/hi/middle_east/8010747.stm}.}

\begin{adjustwidth}{.65cm}{.5cm}
   \begin{singlespace}
    [G]lobal Zionism is the complete symbol of racism, which with unreal reliance on religion has tried to misuse the religious beliefs of some unaware people and hide its ugly face. But what should be seriously considered are the goals of certain superpowers and those in possession of major interests in the world; those who try their best through economic power and political influence and wide media means, to lessen the crimes and ugliness of the nature of the Zionist regime. \ldots. [W]e should try to put an end to the misuse of international means by the Zionists and their supporters. And by respecting nations' demands, we should motivate the united governments to eliminate this clear racism and step on the path of reforming international relations and mechanisms with courage.
    \end{singlespace}
\end{adjustwidth}

\noindent
Such language was a far cry from eight years earlier at Durban I, when the UN High Commissioner for Human Rights at the time (Mary Robinson, former president of Ireland) struck an optimistic note about the opportunities presented by the WCAR, while adding how ``pleased'' she was that ``everyone understands that there can be no return to an issue settled \ldots the former Zionism-racism problem.''\footnote{HR/4554, In Statement To Preparatory Committee Of Racism Conference, Human Rights Commissioner Mary Robinson Reviews Progress, UN Press Release for information media, August 9, 2001. \url{https://press.un.org/en/2001/hr4554.doc.htm}.} 

In fact, that ``former'' problem, never truly laid to rest by UNGA Resolution 46/86 (1991),\footnote{  UNGA Res. 46/86 - English (1991). The Israeli-Palestinian Conflict: An Interactive Database. \url{https://ecf.org.il/media_items/1518}.} which rescinded Resolution 3379, was now completely resurrected, assuming it had ever passed away at all. Unusually entered on the UNGA record without vote in the form of a single conclusory sentence\footnote{G.A. Res. 46/86, U.N. Doc. A/RES/46/86 (Dec. 16, 1991). \url{https://docs.un.org/en/a/res/46/86}.} and appearing to be the only time one UNGA resolution ever revoked another,\footnote{Louis Allday (@Louis\_Allday) (October 10, 2023).[Post] X. \url{https://x.com/Louis_Allday/status/1711789614925496672?s=20}: ``In the entire history of the United Nations, only one successfully ratified General Assembly Resolution has ever been repealed: UNGA 3379 in 1991 after intense US/Israeli pressure and blackmail.'' Interestingly, Mr. Allday, whose articles are often published online in sources such as \textit{Mondoweiss} and \textit{Electronic Intifada}, made this point only three days after October 7, 2023.} this nearly 20-year old step away from UNGA-approved Israel vilification had by then long been forgotten or at most remembered as an expedient tactic, not a change-of-heart reality.\footnote{The 1991 resolution satisfied Israel's insistence that its participation in the Madrid peace conference planned for the  end of that year depended on it, as a quid pro quo to the PLO having observer status, given that Israel had not accepted its purported status at that time as the ``sole legitimate representative'' of the Palestinian people. \textit{See}, UNGA Resolution 46/86 (1991). Economic Cooperation Foundation. \url{https://ecf.org.il/issues/issue/1321}.} It was recognized by informed observers at the time of its passage as a way to remove an impediment to Israel's participation in the Madrid peace conference, jointly organized by the U.S. and newly post-Communist Russia to resolve the ``first intifada'' that began in 1987. That conference succeeded as a prelude to the start of the Oslo process that continued throughout most of the 1990s, through the Rome Statute's ratification and up to the last major agreement in that process (the Wye River Memorandum), both occurring in 1998. Now, nearly ten years later, the Israel--South Africa apartheid analogy emerged in full force, led by the renowned public intellectual and English literature professor Edward Said,\footnote{Edward Said, ``How Do You Spell Apartheid? O-s-l-o,'' October 11, 1998.  Ha'aretz. \url{https://www.edwardsaid.org/articles/how-do-you-spell-apartheid-o-s-l-o/}.} Francis Boyle (a senior advisor to the PLO on international law issues during the Oslo negotiations)\footnote{Boyle, F.A. (Jan 2002) Law and Disorder in the Middle East. The Link, 35, 1. \textit{Americans for Middle East Understanding}, \url{https://www.derechos.org/human-rights/mena/doc/boyle1.html}. (After reviewing the original Oslo agreement, Professor Boyle recalled telling his clients, many years after the fact: ``A bantustan. They are offering you a bantustan. As you know, the Israelis have very close relations with the Afrikaner Apartheid regime in South Africa. It appears that they have studied the bantustan system quite closely. So it is a bantustan that they are offering you.'')} and Norman Finkelstein, a rising Jewish-American academic who had already labeled Israel as ``diabolical''\footnote{Stern-Weiner, J. (March 17, 2014). The end of Palestine? Peace News. \url{https://peacenews.info/node/7563/end-palestine}.} but had not previously compared it to South Africa. At this point, however, he proclaimed that he had always known that the purpose of the Oslo process was to transform Palestine into a ``Bantustan'':\footnote{Finkelstein, N. (Nov. 1998). Securing Occupation: The Real Meaning of the Wye River Memorandum. \textit{New Left Review}. Issue 232. \url{https://newleftreview.org/issues/i232/articles/1977}. Finkelstein's glib mention of South African Bantustans being \enquote{eventually granted statehood} is bizarre, as none of them ever were, except in the most useless, self-serving sense by the political entity that created them in such a way that suggested they should be considered representatively governed \enquote{states.}}

\begin{adjustwidth}{.65cm}{.5cm}
   \begin{singlespace}
    Like the South African Bantustans, the fragmented Palestinian entity resulting from the Oslo process will no doubt eventually be granted statehood. And like the Bantustans, it will be a state in name only \ldots. All the Bantustans won was the right --- in the words of one dissenter --- to `police themselves and administer their own poverty'. This is also the only right Palestinians can expect to win under Oslo.
    \end{singlespace}
\end{adjustwidth}

In such circumstances, it can hardly be seen as surprising that both Durban I and II are now widely remembered not as exercises in careful diplomacy but as public relations disasters in terms of their impact on perceptions of the integrity of the UNGA's humanitarian mission. They were bookends for an all but preposterous decade in the human rights field, when what were presented as lofty steps forward in how best to deal with the scourge of racism were at least cancelled by an equal number of steps back, as UNGA member-states allowed a well-coordinated re-launch of an old national mugging to take place under their noses. A remarkably all-enveloping ``state doctrine'' of racism that had been transparently and officially embraced in South Africa was now repurposed as a house of canards erected solely so that it could be applied to Israel. This was the beginning of a long-term project of finding a way to make orchestrated moral outrage toward Israel ``stick'' even though it was by then a quarter-century old, dating back to Resolution 3379, and apparently was seen as having lost some momentum by virtue of the ongoing 1990s Oslo negotiations. As scholars Noura Erakat and John Reynolds saw it, this was ``mission accomplished'' for both state and non-state actors that had used both conferences expressly for this purpose:\footnote{Erekat, N., \& Reynolds, J. (2021, April 20). We Charge Apartheid? Palestine and the International Criminal Court. \textit{TWAILR: Third World Approaches to International Law Review}. \ \url{https://twailr.com/we-charge-apartheid-palestine-and-the-international-criminal-court/}.}

\begin{adjustwidth}{.65cm}{.5cm}
   \begin{singlespace}
      The controversy which engulfed the UN's 2001 Durban anti-racism conference was precipitated by the insistence of social movements from around the world on calling Israel out as an apartheid state. This was a political and moral argument to challenge institutionalised racism against Palestinians in the context of global anti-racist action. The anti-apartheid framework became increasingly central to political organising and global solidarity through the Palestinian BDS call and initiatives like Israeli Apartheid Week. More recently it has informed renewals of Black-Palestinian solidarities and the Black Lives Matter demand for divestment from apartheid Israel. While Palestinians had diagnosed and detailed Israeli apartheid conditions for many decades prior, the concomitant international legal arguments had only started to be developed in the 1990s by individual scholars and Palestinian rights organisations.
   \end{singlespace}
\end{adjustwidth}

In short order, analogies between 20th century South Africa and 21st century Israel proliferated, both in academia and in grassroots organizing, as exemplified by Stop the Wall, a \enquote{Palestinian grassroots Anti-Apartheid Wall Campaign} that distributed a pamphlet carefully explaining \enquote{why the term apartheid embodies historic, present and future of Palestine \& is a necessary tool for organization and mobilization.}\footnote{This pamphlet, while  no longer found online, laid out the premeditated goal of pushing Israel into pariah state status along the same lines as South Africa by conjuring up a ``public outcry.'' It is referenced in footnote 57 of Keren Greenblatt, ``Gate of the Sun'': Applying Human Rights Law in the Occupied Palestinian Territories in Light of Non-Violent Resistance and Normalization, 12 Nw. J. Int'l Hum. Rts. 152 (2014). \url{https://scholarlycommons.law.northwestern.edu/cgi/viewcontent.cgi?article=1171\&context=njihr}. The organization that published the pamphlet remains active, as shown by its website at \url{https://stopthewall.org/}.} So too did the demands that Israel be treated no differently than South Africa was during its decades of global isolation, through the use of crushing economic, cultural, academic and political boycotts. Thus, in 2005, a loose, single-purpose amalgam of ``Palestinian civil society'' groups --- comprised of NGOs funded mainly from external sources and controlled by political elites with ``almost unlimited powers within their organisations'' \parencite[5]{Shafi-2004}  --- launched \enquote{BDS}, a \enquote{Call for Boycott, Divestment and Sanctions,} describing it as ``inspired by the struggle of South Africans against apartheid.''\footnote{Palestinian Civil Society Call for BDS. (07 Aug. 2005). Palestinian BDS National Committee.  \url{https://bdsmovement.net/call}.} Even Israeli political leaders got the message. In November 2007, after issuing a joint statement with Yasir Arafat at the end of the Annapolis Conference but facing rising dissatisfaction over his leadership at home amidst the violence of the second intifada,\footnote{Palti, Z. (2023, September 8). The Implications of the Second Intifada on Israeli Views of Oslo. The Washington Institute. \url{https://www.washingtoninstitute.org/policy-analysis/implications-second-intifada-israeli-views-oslo}.} Prime Minister Ehud Olmert proclaimed that if the viability of the Oslo process collapsed, ``Israel would face a South African-style struggle for equal voting rights.''\footnote{McCarthy, R. (2007, Nov. 30). Israel risks apartheid-like struggle if two-state solution fails, says Olmert. \textit{The Guardian}. \url{https://www.theguardian.com/world/2007/nov/30/israel}. The news report noted that \enquote{the Israeli government is usually bitterly resentful of any comparison to the apartheid regime, but Olmert's remarks looked like an effort to galvanise support from a skeptical Israeli public for a return to peace negotiations with the Palestinians in the days after the Annapolis conference.}} Peter Beinart, a Jewish-American professor and protector of \enquote{Palestinian freedom} (who also happened to be the son of a South African champion of the anti-apartheid struggle)\footnote{The Beinart Norebook, accessed April 11, 2026. \url{https://peterbeinart.substack.com/}.} joined in as well, writing about how, when Vice President Hillary Clinton \enquote{sped through the West Bank} in 2009 from Jerusalem to Ramallah to meet with Palestinian leaders, one of her aides felt like \enquote{my god, you crossed that border and it was apartheid.}\footnote{Peter Beinart, \textit{The Crisis of Zionism,} Times Books, 2012, 127.}

Within a few years, the visceralness of the analogy gave way to expert assessments that apartheid could indeed be factually established upon careful review of the legal dictates of UNGA conventions and treaties. In 2013, this subject received detailed analysis in a paper co-authored by South African professor John Dugard \parencite{Dugard-2013}, a former UN Special Rapporteur on ``the human rights situation in the Occupied Palestinian Territory'' from 2001 to 2008 and an admirer (along with Noam Chomsky)\footnote{Chomsky, N (2002). The Fate of an Honest Intellectual. \url{https://chomsky.info/power01/}.} of Norman Finkelstein.\footnote{Gaza's Gravediggers - Reviews (Release date: Mar. 31, 2026). \url{https://moonpalacebooks.com/item/c_XdcfkkymNdYa9DHeG9ng}. (``Probably the most serious scholar on the conflict in the Middle East'')} Together with John Reynolds, a professor who by then had ``worked for a number of years in the West Bank for Palestinian human rights organisation Al-Haq,''\footnote{Dr John Reynolds Joins the Department of Law. (30 Jan. 2014). Maynooth University.  \url{https://www.maynoothuniversity.ie/news-events/dr-john-reynolds-joins-department-law}.} they offered a lengthy interpretation of the Rome Statute that authenticated the importance of applying its provisions on apartheid first and foremost to Israel. Thus was born, for all intents and purposes, the Israeli-specific edition of apartheid.

Even so, for at least two reasons, the issue is more complex than the roar of approval from ideology-driven NGOs, or the stamp of authenticity from partisan legal experts, would suggest. First, the use of the phrase ``racial group'' in the Rome Statute  --- as noted earlier, carried over from the ICERD and the Apartheid Convention  --- must be understood as a foundational element in the narrative. Under the Rome Statute's more sophisticated definition of apartheid, the conduct at the heart of the crime must be committed knowingly \enquote{in the context of an institutionalized regime of systematic oppression and domination} but still, this involved what was being done by \enquote{one racial group over any other racial group or groups.}\footnote{Rome Statute, Article 7(2)(h).} This must not have seemed to Dugard and Reynolds as that high of a bar to meet, given that ``racialization'' --- while viewed by some scholars as an inherently subjective analysis --- is meant to be assessed precisely when this sort of practical assessment needs to be considered in a \enquote{historically specific} process \parencite[64]{Omi-Winant-1994}. But as also noted in the prior section, this still requires ``defining precisely what is meant by racialization,'' which at least some more cautious scholars acknowledge \enquote{can be a daunting task, let alone identifying how, when, where, why, and under what conditions it operates.}\footnote{Blanca Gonzalez-Sobrino \&  Devon R. Goss, \enquote{Exploring the mechanisms of racialization beyond the black-white binary,} 2019, Ethnic and Racial Studies, 42.4, 505--510. doi: 10.1080/01419870.2018.1444781. \url{https://doi.org/10.1080/01419870.2018.1444781}.} 

Second, a finding that Palestinians constitute a racial (that is, \enquote{racialized}) group requires further consideration in a location-specific setting which, in this case, must begin with assessing how laws are applied and actions taken by Israelis (both citizens and authorities) in areas deemed by the UNGA (and nearly universally acknowledged) to be \enquote{occupied} within the meaning of the Hague Convention.\footnote{The Law of Belligerent Occupation, 1907 Hague Convention IV, Articles 42--56. \url{https://ihl-databases.icrc.org/ihl/INTRO/195}.}  Much to the consternation of many, some locales (almost all within those identified as Area C in the Oslo Accords) are now interspersed with Jewish communities established with Israeli government approval (pending their final status determination, presumably amicable enough to allow them to continue). These two types of communities --- either Israeli Jewish or Palestinian Muslim (Christian communities are now below 2\% of West Bank residents)\footnote{Giovanna Dell'Orto, \enquote{The West Bank's dwindling Palestinian Christian communities continue to struggle amid violence,} AP New, October 9, 2025 \url{https://apnews.com/article/west-bank-christians-israel-palestinians-mideast-war-912c4f0ea073d8f091e762f1466a7423}.}  --- are obviously entirely distinct. But this too does not require they be viewed as racial groups vis-à-vis one another, much less that their relationship is one of systematic apartheid, because in all areas, the impossibility of intermingled communities reflects both Israeli law and a strong Palestinian preference (through lack of any supporting procedures) not to support any sort of integrated environments. In other words, the fact that these two populations remain segregated in this case, in and of itself, cannot validate the need for a racialization analysis when both sides have, across their long relationship, manifested a clear preference (at least outside what is colloquially known as \enquote{Israel proper}) to maintain separate lifestyles that are fully independent of the other.
 
On a more ad hoc basis, such a state of affairs is not all that uncommon across the MENA region and is often broadly labeled --- although perhaps also unfairly in some cases, because it is simplistic and overly reductive --- as \enquote{sectarianism} \parencite{Makdisi-2017}. This phenomenon may indeed apply to how Israeli settlers interact with Palestinians in what is widely known as \enquote{the West Bank,} and for scholars such as Ussama Makdisi, it may indeed be true that there, rival \enquote{sectarian outlooks, actions, and thoughts \ldots [can be discussed] in a manner similar to how one would talk about racial (and racist) outlooks, actions, and thoughts in the United States} \parencite[4]{Makdisi-2017}. But even if one side of this \enquote{dialectic} garners sympathetic supporters because it is considered \enquote{local} whereas the contrasting side is \enquote{foreign} (as Professor Makdisi has consistently elaborated upon throughout his career), the \enquote{complex, constant, and unequal relationship} \parencite[10]{Makdisi-2017} between them does not always mean that one side is invariably the racist oppressor and the other is racially oppressed. There are, for example, villages within shouting distance of each other in southern Lebanon where residence is based on different religious identities (Christian or Muslim) and, as a result, they have virtually nothing to do with each other, mainly due to long-held hostilities or fears for personal safety, with the Christian villages usually at a power disadvantage relative to the larger Muslim villages.\footnote{E.g., Uzay Bulut, \enquote{From Christian to Majority-Muslim: Lebanon's Cautionary Tale for Europe} (December 25, 2025). The European Conservative. \url{https://www.aljazeera.com/news/2023/4/20/lebanon-clashes-between-christians-and-muslims-in-south}. (\enquote{According to Open Doors, churches in Muslim areas often choose not to display religious symbols and iconography to avoid provoking hostility from the Muslim communities in which they are located. In 2021, clashes between Shia and Christian villages in southern Lebanon's Hezbollah-controlled territory left six people injured. Christian religious symbols were attacked by Shia armed elements during the clashes.})} But such such resulting population zones are seen as having been politically engineered, or are the ineluctable byproduct of ethnic or religion-based spheres of influence, natural preferences for living within what one considers one's own community should not automatically trigger the necessity of undertaking a full-blown racialization analysis.\footnote{Scott A. Bollens, \enquote{Urban Policy in Ethnically Polarized Societies,} 1998, International Political Science Review 19,2, 187--215. \url{https://www.jstor.org/stable/1601323}. While this author is highly critical of Israeli policies in the urban segmentation of Jerusalem, and recognizes that Jewish political dominance disadvantages Arab non-citizen residents there, the issue identified is how best to apply ``practical principles which foster the coexistent viability of antagonistic sides'' (209) in a shared setting, not to assume that distinct populations --- especially those with a long and largely unresolved history of hostility toward each other --- will eventually prefer to live in integrated environments.} 

There is also a third reason why a racial group finding is inappropriate if meant to establish that apartheid is being inflicted upon Palestinians. In South Africa, the system was constructed within a unified state that had unilaterally chosen to demarcate arbitrary \enquote{homelands} or \enquote{townships} --- derisively but accurately labeled as \enquote{bantustans} --- to meet its internal economic and social needs, but without any expectation that they would someday become eligible to obtain any degree of \enquote{self-determination.}\footnote{\enquote{Self-determination} (variously referred to as a \enquote{right} or a \enquote{principle}) is explicitly mentioned in two articles of the United Nations Charter. Its contemporary scope and meaning is generally considered largely derived from Woodrow Wilson's \enquote{Fourteen Points} speech delivered to Congress in January 1918, which did not use the term but included multiple references to the \enquote{political independence and territorial integrity to great and small states alike.} As noted in a 1965 law article, The \enquote{essence} of Wilson's approach was to apply the notion to the \enquote{self-government of peoples.} M.K. Nawaz, {The Meaning and Range of the Principle of Self-Determination,} Duke Law Journal, Vol. 1965, No. 1, pp. 82--101. \url{https://scholarship.law.duke.edu/cgi/viewcontent.cgi?article=1918\&context=dlj}. However, the interest of the UNGA to advance such \enquote{peoples} related to its collective desire to bring about \enquote{a speedy and unconditional end [to] colonialism in all its forms and manifestations.} UNGA Resolution 1514 (XV) (December 14, 1960), \url{https://www.ohchr.org/en/instruments-mechanisms/instruments/declaration-granting-independence-colonial-countries-and-peoples} Clearly, the South African bantustan system was the antithesis of this desire, as it represented fragmenting peoples into artificially localized areas of supposed self-control that had no connection to meaningful sovereignty. This is clear from an interview that Prime Minister P.W. Botha gave to Time Magazine's local bureau chief in 1979. \enquote{Putting a Pretty Face on Apartheid,} Time Magazine, December 3, 1979, 59--60. Referring to the country's \enquote{native} population, he noted that \enquote{we have made available \ldots certain facilities in housing, which is welcomed by most of them}; and \enquote{we will grant them local government so they will be able to control many aspects of their lives.} The opportunity for eventual self-determination and political independence was never anticipated, much unlike the intention reflected by the ``interim self-government'' wording used in the 1993 Oslo accords.} But unlike the situation in the Palestinian territories, enforced limitations on population movement within South Africa clearly took place inside a unitary nation-state entity, which unilaterally pretended to grant some of its inhabitants \enquote{local} rights for purposes of precluding their exercise of rights more broadly. Such a situation could not be any further removed from the UNGA's continual involvement in the creation of a standalone and fully independent \enquote{State of Palestine} since its inception, with the ongoing involvement and support of national governments around the world.  Juxtaposing that with the frustrations felt by Palestinians for their lack of such a state, even if they fairly believe its natural development has been willfully stunted and suppressed by Israelis, one difference should be abundantly clear: finding a path toward their eventual national autonomy is not the problem that prosecuting Israelis for apartheid crimes is meant to solve. That is, assuming that they may still achieve a separate nationality (and in some respects they already have, given Palestine's recognition as a state by many other states), their ultimate success is not dependent on addressing the purportedly racist (and therefore internal) actions of citizens of another state but instead resuscitating a process toward achieving a viable diplomatic (and therefore external) relationship with that state. 

Taking another angle on why tensions between Israelis and Palestinians cannot be understood as a purely \enquote{racial} issue, another highly impactful factor that has plagued the continuously troubled status of the Palestinian territories relates to the looting of the public funds of the Palestinian Authority which, together with a ``bipolar system'' that splintered and factionalized its political alliances --- even while all purport to \enquote{represent democratic and secular elements [within the] society} \parencite[10]{Shafi-2004} --- has made it impossible for nascent public institutions to deliver essential services \parencite {Samuels-2005}. This phenomenon has continued since the Palestinian Authority's inception and notwithstanding tremendous inflows of international aid, has had state-formation crushing implications.\footnote{As noted by Samuels, \enquote{[t]he amounts of money stolen \ldots through the corrupt practices of Arafat's inner circle are so staggeringly large that they may exceed one half of the total of \$7 billion in foreign aid contributed to the Palestinian Authority. \ldots The International Monetary Fund has conservatively estimated that from 1995 to 2000 Arafat diverted \$900 million from Palestinian Authority coffers, an amount that did not include the money that he and his family siphoned off through such secondary means as no-bid contracts, kickbacks, and rake-offs. A secret report prepared by an official Palestinian Authority committee headed by Arafat's cousin concluded that in 1996 alone, \$326 million, or 43 percent of the state budget, had been embezzled, and that another \$94 million, or 12.5 percent \ldots went to the president's office, where it was spent at Arafat's personal discretion. An additional 35 percent of the budget went to pay for the security services, leaving a total of \$73 million, or 9.5 percent of the budget, to be spent on the needs of the population of the West Bank and Gaza.}} Nonetheless, the duplicity of Israel, due to its racism, is imputed to be the sole cause of Oslo's eventual failure irrespective of the stunning levels of corruption of Palestinian leaders and public figures who have failed to uphold their end of the \enquote{two-state solution} bargain: developing a long-term process for building a viable state of their own. A more balanced assessment of this outcome, and the failed \enquote{risk} taken on the Palestinian side which ultimately enervated its chances for long-term positive results, was noted in passing by a Harvard psychology professor in 1997 \parencite[194]{Kelman-1997}:

\begin{adjustwidth}{.65cm}{.5cm}
   \begin{singlespace}
    For the Palestinians, an agreement that did not guarantee an independent Palestinian state at the end of the interim period \ldots was highly risky. They had to place a great deal of trust in the dynamics of the negotiation process, which Arafat seemed prepared to do. What made it possible for him to accept the risk were two features of the Oslo accord: the territorial base and the early empowerment of the Palestinian Authority \ldots. Early empowerment assured the Palestinian leadership that they could immediately begin the process of \textit{establishing control and building pre-state institutions} (emphasis added).   
  \end{singlespace}
\end{adjustwidth}

This may seem a tangential consideration nearly thirty years after the hopes for the Oslo process began fading and more than twenty years after Yasser Arafat's death, but the problem has been ongoing and reveals the lacuna between Palestine's fragmented and fragile political reality (especially the dichotomy between \enquote{Islamist} and \enquote{secular} power bases)\footnote{Add a footnote here.} and whatever on-the-ground conclusions might be reached by taking a look around and noticing a marked divergence between the relative levels of prosperity between Jewish and Arab communities in the territories. Notwithstanding the gut reactions (as reported by Peter Beinart) of a member of Hillary Clinton's entourage in 2009, this divergence cannot automatically trigger applying a paradigm of one group having undergone a racialization process due to the presence of the supremacist \enquote{Other}, thereby making it susceptible to oppression. True, this concept has now been deeply explored by sociologists and political theorists examining the basis of racial tensions in diverse national societies. But their analytical model deals with foregrounding \enquote{the important justice-based goal of keeping issues of social inequality in the forefront of our thinking about such groups} \parencite[304]{Blum-2010}, not on failed co-existence negotiations that have soured even further due to the constant exacerbation of internal tensions, whether due to ideological chasms, random violence or purposefully violent \enquote{resistance} or \enquote{failure-to-launch} state institution building. 

Turning to the international law aspects of this overall situation, also of possible relevance is a decided case involving the scope of authority of the Committee on the Elimination of Racial Discrimination (CERD). The CERD is empowered by section 8 of the ICERD to handle cases between member states brought before the International Court of Justice (ICJ) relating to racially discriminatory treatment of residents. But in a 2021 case between Qatar and the UAE, the ICJ rejected a position taken by the CERD and ruled that the ICERD did not extend to claims made by Qatar alleging \enquote{nationality-based discrimination.}\footnote{Qatar v. United Arab Emirates, Judgment of February 4, 2021, ICJ Reports 2021, 
\url{https://www.icj-cij.org/public/files/case-related/172/172-20210204-JUD-01-00-EN.pdf}. The ICJ ruling clarified that the ICERD's inclusion of the term \enquote{national origin} as a basis for racial discrimination claims related to the protection of \enquote{persons of foreign origin who are subject to racial discrimination in their country of residence on the grounds of that origin} (98, \P 51.)} Therefore, a \enquote{racial group} finding cannot be bypassed simply because Israelis are residing in areas considered by the UNGA to be \enquote{occupied Palestinian Territories} and are interacting in allegedly illegal ways with all other residents who are not Israelis. This ruling is consistent with a 2007 ICJ ruling involving the Genocide Convention, which is similar in its scope of protection to the ICERD in its requirement that claims can only be brought on behalf of a specified group of individuals consisting of a substantial number of that group (ICJ 2007 Ruling, 124, \P 193).\footnote{International Court of Justice (February 26, 2007). Case Concerning Application of the Convention on the Prevention and Punishment Of The Crime Of Genocide --- Bosnia and Herzegovina v. Serbia and Montenegro, 124. \url{https://www.icj-cij.org/sites/default/files/case-related/91/091-20070226-JUD-01-00-EN.pdf}. (\enquote{The intent must also relate to the group \enquote{as such}. That means that the crime requires an intent to destroy a collection of people who have a particular group identity.})} Moreover, the ICJ has also noted that meeting the \enquote{group requirement} does not always depend on the \enquote{substantiality requirement} alone, although it is \enquote{an essential starting point} (ICJ 2007 Ruling, 127 \P 200). And yet, rather than assess whether Palestinian Arabs all belong to one all-encompassing ``racial group'' or whether a substantial number within that group is actually being subjected to apartheid, these definitional prerequisites are routinely glossed based solely on a presumed oppositional relationship with all Israelis, even by the ICJ's current president, Lebanese former politician, diplomat and jurist Nawaf Salam.\footnote{Declaration of President Salam (July 19, 2023), Document Number 186-20240719-ADV-01-01-EN in Case 186 (Legal Consequences arising from the Policies and Practices of Israel in the Occupied Palestinian Territory, including East Jerusalem), 5--6 (paras 21-22). \url{https://www.icj-cij.org/index.php/node/204161} Para. 21 (``Article 1 of [ICERD] clearly indicates that \enquote*{race} is not the only criterion of racial discrimination. \ldots Palestinians and Israeli Jews identify themselves as two distinct groups based on \enquote*{subjective} elements relating to \enquote*{descent, or national or ethnic origin}, including those relating to religion and culture, and should therefore be considered as two \enquote*{racial groups} within the meaning of the first constitutive element of apartheid.'')} 

In re-examining racialization as the process by which a racial group is formed, it is also noteworthy how widely agreed upon the concept is in academic circles as based on the ground-breaking, yet also in some ways strikingly ambiguous, expostulation of its core features by American anthropologists Michael Omi and Howard Winant in the early 1990s.\footcite[][pp. 53-54. (\enquote{A second approach has been to focus on the common processes by which groups are formed (and reformed) as racial groups—that is, are identified by social and political institutions and self-identify as distinct races. This approach bypasses the problem of mapping racialized groups in a conceptual space or in a hierarchy of groups. Michael Omi and Howard Winant took this approach in their seminal work, Racial Formation in America, originally published in 1989 and reissued in revised versions in 1994 and 2014.})]{Glenn-2015} In their book \textit{Racial Formation in America,} still widely taught after its initial publication in 1989, these two scholars conceived of it as a segmentation process by which a sub-population within a specific social organization that is national in scope (in their case, the United States) becomes \enquote{racialized} by the actions of more dominant groups exercising economic, social and political power over it. This is hypothesized to result in the sub-population's subjugation within a racial context, but even their analysis seems uncertain of how and when a racialized consciousness emerges or what exactly its ramifications are when manifested \parencite[138]{Omi-Winant-1994}:

\begin{adjustwidth}{.65cm}{.5cm}
   \begin{singlespace}
   We have approached race as a phenomenon whose meaning is contested throughout social life. Race is a constituent of the individual psyche and of relationships among individuals; it is also an irreducible component of collective identities and social structures. In American history, racial dynamics have been a traditional source both of conflict and division and of renewal and cultural awareness. Race has been a key determinant of mass movements, state policy, and even foreign policy in the United States. Racial meaning systems are contested and racial ideologies mobilized in political relationships.
   \end{singlespace}
\end{adjustwidth}

Nonetheless, this approach treats the overall process of racial group formation --- which may or may not be tied to conventional understandings of race as based on obvious physical differentiators between oppressed and oppressor groups ---  as \enquote{a fundamental organizing principle of social relationships} within the country where it takes place \parencite[66]{Omi-Winant-1994}. Any inquiry into whether the process has in fact occurred in a particular case necessitates a location- and population-specific situational analysis that presumably cannot extend to any and all populations or social or political power dynamics within that country. Rather, according to Omi and Winant's thinking, racial differences are best understood as ``a fundamental axis of social organization \ldots that cannot be subsumed within a supposedly more fundamental category'' that is not biological in origin \parencite[13]{Omi-Winant-1994}, thereby accentuating that both the \enquote{fundamentalness} of the resulting racial group, and the singularity of the ``social organization'' in which its racialization process unfolded, must be considered separate aspects that are both essential to the analysis. 

This thinking is not altogether easy to parse, especially because it is offered as a byproduct of how the formerly \enquote{dominant paradigm of ethnicity, the centerpiece of liberal racial politics, fell into disrepute among radicals} because of the perceived insufficiency of the \enquote{explanatory power and political efficacy of the ethnicity perspective} \parencite[139]{Omi-Winant-1994}. But what should be clear is that before a court decides to undertake an inquiry into the presence of a racial group, as it chose \textit{not} to do in the Qatar vs. UAE case,  the ICERD's provisions barring racial discrimination and apartheid, as well as the Rome Statute's criminalization of apartheid, can only be considered if brought by or on behalf of individuals (not states) belonging to the impacted racial group. This is especially true because as a crime against humanity, the Rome Statute can only be used to prosecute against individuals.\footnote{International Criminal Court --- About the Court (n.d.) \url{https://www.icc-cpi.int/about/the-court}. (\enquote{The International Criminal Court (ICC) investigates and, where warranted, tries individuals charged with the gravest crimes of concern to the international community: genocide, war crimes, crimes against humanity and the crime of aggression.})} The issue then becomes what enables a racialization analysis relating to all Israelis against all Palestinians based on consideration of them as separate racial groups as opposed to separate nationalities? 

Support for the conclusion that such an analysis is not possible can be derived from their mutual desire not to share citizenship in a single \enquote{binational} state.\footnote{Ghada Karmi, ``The One-State Solution: an Alternative Vision for Israeli-Palestinian Peace.'' \textit{Journal of Palestine Studies}, Vol. XL, 2 (Winter 2011), 62--76, 72. \url{https://www.palestine-studies.org/sites/default/files/attachments/jps-articles/jps_2011_xl_2_62.pdf}. (``The reality is that opponents of the one-state solution far outnumber supporters. Much of the resistance --- from both Palestinians and Israelis --— relates to the formidable obstacles to the idea's implementation, starting with the lack of constituency on either side. Indeed, a 2009 Jerusalem Media and Communication Centre opinion poll found that only 20 percent of West Bank and Gaza Palestinians, and just 9 percent of Jewish Israelis favored a binational solution.'')} And yet, as ICJ President Salam suggested in his 2024 declaration\footnote{Mark Goldfeder (@MarkGoldfeder) (July 19, 2024).[Post] X. \url{https://x.com/Louis_Allday/status/1711789614925496672?s=20}: \enquote{Remember that the President of the Court, Judge Nawaf Salam, previously served as Lebanon's ambassador to the United Nations. During his tenure as ambassador Salam routinely accused Israel of apartheid, war crimes, and terror, and sided with Iran. Why is this important? Because under the ICJ's charter (Article 17, Section 2) \enquote{No member may participate in the decision of any case in which he has previously taken part as \ldots advocate for one of the parties \ldots or in any other capacity.}}} and professors Reynolds and Dugard concluded in their 2013 paper, because ``Jews and Palestinians constitute distinct `racial groups' for definitional purposes,'' a finding of apartheid is possible and need not be limited to ``a narrow construction of race,'' with that ``crucial preliminary question [having thereby] been met'' \parencite[885-889]{Dugard-2013}. In effect, their relationship could thus be considered equivalent to a social utilization of the concept of race \enquote{to organize and distort perceptions of the world's various populations,} even though their differences did \enquote{not sit within traditional conceptions of \enquote*{race}} \parencite[887]{Dugard-2013} and arose entirely outside a purely domestic (that is, intra-national) setting.

In so doing, this jurist and these professors have collapsed all differences between ``the world's various populations'' into an exercise that can ascertain a state of discrimination between them, even though they have no shared obligations as members of the same national society. This turns any human rights inquiry relating to racism and racialization on its head, changing it from societal situations relating to a country's sub-populations into an issue with explosively broad, international ramifications. If it is possible to criminalize on racial grounds the existence of any conflict between nations whenever one side is deemed stronger and the weaker side alleges it is forced to act against its will because it was \enquote{racialized} by the stronger side, what need is there for laws pertaining to the conduct of war or international relations, as the stronger side will always be in the wrong and the weaker in the right? This could have been what the ICERD, the Apartheid Convention, or the Rome Statute intended and the very fact that no dispute brought on such grounds has ever been submitted to the ICJ or ICC for briefing reflects that it was not. In fact, the WCAR Declaration and Programme of Action from Durban I, addressing ``Victims of Racism, Racial Discrimination, Xenophobia and Related Intolerance,'' makes this clear. In its sole use of the phrase ``racial group,'' it states that all UNGA member-states must recognize ``the effect that discrimination, marginalization and social exclusion have had and continue to have on many \textit{racial groups living in a numerically based minority situation within a State},'' noting that all such groups should enjoy ``all human rights and fundamental freedoms without distinction and in \textit{full equality before the [presumably singular] law \ldots  [with] respect of employment, housing and education with a view to preventing racial discrimination.''}\footnote{Declaration and Programme of Action (2002). World Conference Against Racism, Racial Discrimination, Xenophobia and Related Intolerance. UN Department of Information, New York. \url{https://www.ohchr.org/sites/default/files/Documents/Publications/Durban_text_en.pdf}.} 

A final related issue is whether Israeli policy towards Palestinian communities outside its borders is contrary to the ICERD's admonition against ``racial segregation,'' which requires fair and equal treatment of a country's inhabitants. This issue is covered above but only as an indicator of whether a racial group can be found to exist. Outside of that, it that has not been codified as a ``crime against humanity'' or even specifically recognized as a crime under international law. But in a 2023 case in which the UNGA requested an advisory opinion from the ICJ,\footnote{International Court of Justice (19 Jul. 2024). Legal Consequences Arising from the Policies and Practices of Israel in the Occupied Palestinian Territory, Including East Jerusalem. \url{https://www.icj-cij.org/sites/default/files/case-related/186/186-20240719-adv-01-00-en.pdf}. p. 7.} the issue presented was whether Israel had violated Palestinian human rights by enacting ``legislation and measures [that] \ldots maintain a near-complete separation in the West Bank and East Jerusalem between the settler and Palestinian communities.'' The ICJ advised that it had, referring to the ICERD's condemnation of ``racial segregation and apartheid'' in Article 3. But it did not specify whether its finding was based on both or either of these prohibitions, leading to interesting ``differences of views on the bench.''\footnote{Victor Kattan, ``Apartheid or Systemic Discrimination? A Connotative Reading of the ICJ's Advisory Opinion,'' Oct. 17, 2024, \textit{Verfassungsblog}, \url{https://verfassungsblog.de/apartheid-or-systemic-discrimination/}.}  Rather than assert, as has one scholar who is sympathetic to Palestinian concerns has done, that the ICJ ``made an implicit apartheid finding going beyond segregation,''\footnote{Id.} it seems more likely that the ICJ chose explicitly not to address that issue at all. Instead, it appears to have found such separation is reflective of \enquote{racial discrimination} because the concept's breadth covers any \enquote{distinction, exclusion, restriction or preference based on race, colour, descent, or national or ethnic origin} (ICERD, Art. 1, sec. 1). It could also have fairly assumed that it had not been asked to opine on whether apartheid --- as a crime against humanity --- had  occurred because any such finding would have depended on a ``subjective element'' to commit the crime that was not factually in the record before it and therefore beyond its purview.\footnote{This point is noted in the separate opinion of Judge Nolte published alongside the court's final opinion. Separate Opinion of Judge Nolte, \url{https://www.icj-cij.org/sites/default/files/case-related/186/186-20240719-adv-01-08-en.pdf}. \textit{Dolus specialis} --- that is, a \enquote{specific intent to cause a specific kind of harm} --- is required whenever a crime is deemed heinous, based on the widely accepted legal principle that \enquote{specific intention, required as a constitutive element of the crime}  depends on showing that \enquote{the perpetrator clearly [sought] to produce the act charged.} Guide to Latin in International Law, 2d ed., 83.}

In sum, given the partisanship of several of the scholars who have examined the subject, their views --- in and of themselves --- should be seen at most as the instantiation of a ``3.0'' iteration of the crime of apartheid, with the ``1.0'' version consisting of apartheid as practiced by South Africa before 1994 and the ``2.0'' version consisting of its internationalization as a crime under the Apartheid convention. This latest updated version appears mostly designed as a receptacle for the outrage generated over the universal acclaim that a \enquote{crime against humanity} has in fact been committed, not the allegations made in a legal case that would lay out the supporting facts to prove the grounds for such outrage. The goal seems to be to continue buttressing what Israel's many haters see as ``growing international recognition of Israel's crime of apartheid against the Palestinian people.''\footnote{Cairo Institute for Human Rights Studies (16 Mar. 2022) Coinciding with the 49th UN Human Rights Council: 16-day awareness raising campaign on Israel's crime of apartheid against the Palestinian people. \url{https://cihrs.org/coinciding-with-the-49th-un-human-rights-council-16-day-awareness-raising-campaign-on-israels-crime-of-apartheid-against-the-palestinian-people/?lang=en}.} But this ultimately becomes an echo chamber that exists in the absence of either impartial trials or clear judicial outcomes. As such, the consensus views of self-proclaimed experts that \enquote{Israel is an APARTHEID STATE}\footnote{\enquote{Israel is an APARTHEID STATE --- The B'Tselem Analysis,} A briefing paper by the Ireland-Palestine Solidarity Campaign - March 2021. \url{https://www.ipsc.ie/wp-content/uploads/2021/07/Apartheid-The-BTselem-Analysis-IPSC-Briefing-March-2021.pdf}.} must be understood as bearing the imprint of those who have long held intractably antagonistic views toward Israel. In the words of an Egyptian institute dedicated to the study of human rights, its purpose may be directed primarily toward reinforcing societal assent to the hostile talking points of its Arab state sponsor:\footnote{Id.}

\begin{adjustwidth}{.65cm}{.5cm}
   \begin{singlespace}
    This growing international and UN recognition would not have been possible without the pressure and advocacy campaigns of Palestinian, regional, and international civil society groups and human rights organizations. These campaigns have brought global awareness to how Israel imposes an apartheid regime on the Palestinian people through an arsenal of repressive and racist laws, policies, practices, including violence and intimidation.   
  \end{singlespace}
\end{adjustwidth}

In the realm of legal adjudications by the court of last resort within a sovereign nation, it may be that determining something to be true effectively makes it so, at least in the sense that the determination has the effect of becoming part of the binding law within that nation. But in the case of international law, which exists within spaces created between nations through bodies considered authoritative in their shared framework of international relations, ICJ advisory opinions are not the same as rulings based on findings of proven facts. Thus, such opinions are just that --- not legally enforceable --- and as pronouncements of ``global awareness'' of the validity of the opinions, they are not always accurately depicted when described as ``global'' or ``awareness.'' Regardless, as there clearly is widespread support for the view that Palestinians must be seen as enduring steadfastly against the racist oppressions of Israelis, the antecedents of this embedded frame of mind deserves deeper scrutiny, which are now taken up in the following section. 
\parindent 0pt
\parskip 6pt

\setenotez{
  list-name = {Notes to Chapter \thechapter}
}
% Place this macro at the exact end of every chapter:
{\footnotesize
\printendnotes[paragraph]}

\thispagestyle{empty} % makes a blank page

\end{document}